% !TeX program = xelatex
% Adapted from https://github.com/emslop/laser-welding-biotissues/blob/main/presentation.tex
% Source-only template: no PDF build is required.
\documentclass[aspectratio=169,11pt]{beamer}
\usepackage{fontspec}
\setsansfont{Liberation Sans}
\setmainfont{Liberation Sans}
\setmonofont{Liberation Mono}
\usepackage{polyglossia}
\setdefaultlanguage{russian}
\usepackage{amsmath,booktabs,tabularx,array,graphicx,tikz,pgfplots,pgfpages,etoolbox}
\usetikzlibrary{arrows.meta,positioning,calc}
\pgfplotsset{compat=1.18}
\graphicspath{{./}{images/}{figures/}}
\usefonttheme{professionalfonts}
\setbeamertemplate{navigation symbols}{}
\setbeamertemplate{headline}{}
\setbeamertemplate{frametitle}{\vspace{2mm}\centering\bfseries\insertframetitle\par}
\setbeamercolor{normal text}{fg=black,bg=white}
\setbeamercolor{frametitle}{fg=black,bg=white}
\setbeamercolor{itemize item}{fg=black}
\setbeamercolor{enumerate item}{fg=black}
\setbeamertemplate{itemize item}{\textemdash}
\setbeamertemplate{itemize subitem}{\textemdash}
\setbeamertemplate{footline}{\hfill\insertframenumber/\inserttotalframenumber\hspace{5mm}\vspace{3mm}}
\setbeamersize{text margin left=7mm,text margin right=7mm}
\setlength{\parskip}{3pt}
\BeforeBeginEnvironment{frame}{\setcounter{figure}{0}\setcounter{table}{0}\setcounter{equation}{0}}
\newcommand{\figcap}[1]{\par\refstepcounter{figure}{\centering\normalfont\mdseries\upshape\fontsize{9}{10.8}\selectfont Рисунок~\thefigure~--- #1\par}\vspace{1mm}}
\newcommand{\tabcap}[1]{\par\refstepcounter{table}{\centering\normalfont\mdseries\upshape\fontsize{9}{10.8}\selectfont Таблица~\thetable~--- #1\par}\vspace{2mm}}
\newcommand{\pngimage}[2]{\includegraphics[width=\linewidth,height=#2,keepaspectratio]{#1}}
\begin{document}
\begin{frame}[plain]
\begin{tikzpicture}[remember picture,overlay]
\node[anchor=north,inner sep=0pt] at ([yshift=-3.3mm]current page.north) {\parbox{\dimexpr\paperwidth-12mm\relax}{\centering\fontsize{8.5}{10.2}\selectfont\setlength{\parskip}{0pt}%
\mbox{МИНОБРНАУКИ РОССИИ}\\[2pt]
\mbox{Федеральное государственное автономное образовательное учреждение высшего образования}\\
\mbox{«Национальный исследовательский университет}\\
\mbox{«Московский институт электронной техники»}\par}};
\node[inner sep=0pt] at ([yshift=-20.8mm]current page.north) {{\fontsize{8.5}{10.2}\selectfont Институт биомедицинских систем}};
\node[inner sep=0pt] at ([yshift=-33.8mm]current page.north) {{\fontsize{14}{17}\selectfont Иванов Иван Иванович}};
\node[inner sep=0pt] at ([yshift=-52mm]current page.north) {\parbox{\dimexpr\paperwidth-12mm\relax}{\centering\fontsize{15}{17.5}\selectfont\bfseries\setlength{\parskip}{0pt}%
Шаблон презентации\par}};
\node[anchor=east,inner sep=0pt] at ([xshift=-9mm,yshift=13.8mm]current page.south east) {{\fontsize{8.5}{10.2}\selectfont Руководитель: к.ф.-м.н., доцент Потапов Д.А.}};
\node[inner sep=0pt] at ([yshift=4mm]current page.south) {{\fontsize{8.5}{10.2}\selectfont Москва 2026}};
\end{tikzpicture}
\end{frame}

\begin{frame}{Структура системы с обратной связью}
\begin{columns}[T,onlytextwidth]
\begin{column}{0.55\linewidth}
\centering\pngimage{example-system.png}{35mm}
\par\vspace{1mm}
{\fontsize{8.5}{10.2}\selectfont
1 — источник энергии; 2 — преобразователь;\\
3 — нагрузка; 4 — датчик; 5 — блок управления\par}
\figcap{Условная система с обратной связью}
\end{column}
\begin{column}{0.42\linewidth}
\tabcap{Назначение компонентов системы}
{\fontsize{10}{12}\selectfont
\begin{tabularx}{\linewidth}{@{}>{\centering\arraybackslash}p{6mm}>{\raggedright\arraybackslash}X@{}}
\toprule
\textbf{№} & \textbf{Назначение} \\
\midrule
1 & Подача энергии. \\
2 & Преобразование энергии. \\
3 & Выполнение полезной работы. \\
4 & Измерение выходной величины. \\
5 & Формирование управляющего сигнала. \\
\bottomrule
\end{tabularx}}
\end{column}
\end{columns}
\note{Учебный пример оформления, а не схема конкретного прибора и не экспериментальные данные. Верхняя цепь показывает передачу энергии; нижняя ветвь — измерение и управление.}
\end{frame}
\end{document}
